% Separate post-CP27 proved corollary. Four complete audits passed; closure of the same bounded-period hereditary projective enlargement.
\subsection{Bounded-period projective crowns over arbitrary permutation quotients}

\begin{corollary}[Proved: arbitrary-quotient projective crowns with bounded multiplier period]
\label{D:bounded-period-arbitrary-quotient-projective-crowns}
Let $R$ and $\Gamma$ be as in
Theorem~\ref{D:bounded-period-hereditary-projective-growth}.
Fix constants
$B_*,b_0\ge1$ and $\kappa>0$.
Let $L$ be a finite monolithic primitive group with
nonabelian socle $N=S^t$, where $S$ is a finite nonabelian
simple group, and put $P=L/N$.
Fix a continuous epimorphism $\Gamma\to P$.
Choose a socle factor, and let $\mathcal B\le\operatorname{Aut}(S)$
be the image of its stabilizer in $L$. Suppose
\[
 S\le\mathcal B_0\le\mathcal B,
 \quad [\mathcal B:\mathcal B_0]\le b_0,
 \quad \iota:\mathcal B_0\hookrightarrow
                     \operatorname{PGL}_r(\mathbb F_q),
 \quad r^2\log q\le\kappa\log|S|,                    \tag{D.BC.1}
\]
where $q$ is a power of the defining characteristic $p$.
Identify $S$ with its inner automorphism group.

Let $H\subset\Gamma$ be the preimage of the point
stabilizer of that factor. The fixed outer map determines
$\eta:H\to\mathcal B/S$, and put
$H_0=\eta^{-1}(\mathcal B_0/S)$.
Require that, for every stabilizer lift $\theta:H\to\mathcal B$
of $\eta$ which occurs in the crown comparison, the
projective representation
\[
                  \iota\theta|_{H_0}:H_0\to
                                  \operatorname{PGL}_r(\mathbb F_q)
                                                        \tag{D.BC.2}
\]
has geometric Brauer period at most $B_*$ on the puncture
corresponding to $H_0$. This condition holds for all such
lifts if $\gcd(r,q-1)\le B_*$.

If the crown-based power $L_k$ is an epimorphic image
of $\Gamma$ over the chosen map to $P$, then
\[
 k\le |N|^{c_R},\qquad
 d(L_k)\le\max\{2,d(P),c_R+2\},                       \tag{D.BC.3}
\]
where one may take
\[
 c_R=\max\left\{1,
   \left\lceil\kappa(A_{R,B_*}b_0+1)+b_0-1\right\rceil\right\}.
                                                        \tag{D.BC.4}
\]
There is no tameness or etale-extension assumption on
the permutation quotient cover. The term $d(P)$ is
retained. The embedding in (D.BC.1) is an ordinary
projective embedding; a semilinear embedding does not
satisfy this hypothesis by itself.
\end{corollary}

\begin{proof}
The exact nonsplit transport proof of
\ref{D:nonsplit-automorphism-wreath-lifts}, and the
restriction argument of
\ref{D:normal-subgroup-projective-crowns}, identify the
relevant $N$-conjugacy classes of lifts over $P$ with
$S$-conjugacy classes of stabilizer lifts
$\theta:H\to\mathcal B$ having the prescribed outer map
$\eta$. We recall the points needed for this use.
The centralizer of the unique nonabelian socle is trivial,
so conjugation gives the faithful actual automorphism
wreath embedding. Choose actual coset transports between
the socle factors. Restriction to the chosen factor gives
the stabilizer lift, and changing the transports changes
it by inner $S$-conjugation. The reverse construction
uses the same actual transports and does not split the
outer extension. Two crown coordinates in the same
$N$-conjugacy class would force their pair to lie in a
diagonal graph subgroup, contradicting surjectivity onto
the crown-based power. Therefore $k$ is at most the
number of these stabilizer lift classes.

Put $a=[H:H_0]\le b_0$. The open subgroup $H$ has
index $t$ in $\Gamma$, and $H_0$ has index $at$.
Every lift maps $H_0$ into $\mathcal B_0$, by its fixed
outer map. The assumed multiplier-period condition and
the hereditary projective theorem give at most
\[
                         q^{A_{R,B_*}at r^2}            \tag{D.BC.5}
\]
ambient $\operatorname{PGL}_r(\mathbb F_q)$ conjugacy
classes of the restricted maps. Faithfulness of $\iota$
identifies an actual restricted matrix homomorphism with
at most one homomorphism to $\mathcal B_0$.
Each ambient orbit contains at most
$|\operatorname{PGL}_r(\mathbb F_q)|\le q^{r^2}$
actual homomorphisms, hence no more than that many
$S$-conjugacy classes. This uses the ambient orbit
cardinality, without any normalizer-conjugacy assertion.
Thus the restricted $S$-class count is at most
\[
       q^{(A_{R,B_*}at+1)r^2}
                   \le |S|^{\kappa(A_{R,B_*}at+1)}.     \tag{D.BC.6}
\]

For each restricted $S$-class choose an actual
representative. Conjugate a full lift by $S$ so that
its restriction is this representative. Choose $a-1$
nonidentity coset representatives for $H_0$ in $H$.
Together with $H_0$ they determine every homomorphism
from $H$. The fixed outer map places the image of
each such representative in one fixed coset of $S$
in $\mathcal B$, giving at most $|S|$ choices each.
Consequently extension costs at most $|S|^{a-1}$.
This argument needs neither normality of $H_0$ nor
normality of $\mathcal B_0$; $S$ itself is normal
in every group displayed in (D.BC.1).
Combining with (D.BC.6) gives
\[
 \#\{\theta\}/S
   \le |S|^{\kappa(A_{R,B_*}at+1)+a-1}
   \le |S|^{t\{\kappa(A_{R,B_*}b_0+1)+b_0-1\}}
   \le |N|^{c_R}.                                    \tag{D.BC.7}
\]
The exact crown generation estimate already used in
the cited transport and normal-subgroup corollaries,
from the fully audited finite-group input, now gives
the second inequality of (D.BC.3). Its hypotheses are
the nonabelian simple socle, the bound $k\le|N|^{c_R}$,
and the actual quotient $P$; no bound for $d(P)$ is
inserted.

Finally, the determinant of chosen linear lifts of
a projective representation makes its full scalar
obstruction $r$-torsion, while its coefficient group
has exponent $q-1$. Its geometric image therefore
has period dividing $\gcd(r,q-1)$, proving the stated
sufficient condition.
\end{proof}

\begin{remark}[A precise central-lift alternative]
Another sufficient condition in (D.BC.2) is the following
explicit one. Suppose the pullback to $\mathcal B_0$ of
the scalar extension
$\operatorname{GL}_r(\mathbb F_q)\to
\operatorname{PGL}_r(\mathbb F_q)$ is the pushout of
a fixed central extension
$1\to C\to\widetilde{\mathcal B}_0\to\mathcal B_0\to1$
with $\exp C\le B_*$.
Every multiplier in (D.BC.2) is then the pushforward
of a class in $H^2(H_0,C)$ and is killed by $\exp C$.
Its geometric period is at most $B_*$.
Merely naming a small abstract centre without this
actual compatible central lift would not prove the
required coefficient-class assertion.
\end{remark}

\begin{remark}[A verified ordinary type-A instance]
The natural projective model for the simple groups
$S=\operatorname{PSL}_r(\mathbb F_q)$ is among the
fully checked defining-characteristic models of
\ref{D:classical-defining-characteristic-crowns}.
Apart from the stated nonsimple small cases there,
it has the uniform order comparison with $\kappa=16$.
Thus one may take any
$S\le\mathcal B_0\le\operatorname{PGL}_r(\mathbb F_q)$
and any stabilizer image $\mathcal B\le\operatorname{Aut}(S)$
containing it with $[\mathcal B:\mathcal B_0]\le b_0$,
provided $\gcd(r,q-1)\le B_*$.
The natural inclusion is the actual faithful projective
embedding; the cover is allowed to be arbitrary.
No recognition of an unbounded field-automorphism image
inside an ordinary projective group is being asserted.
This is one concrete scope of the same bounded-period
enlargement, not a second counting reduction.
\end{remark}
